\documentclass[11pt,article,oneside]{memoir}
\usepackage[utf8]{inputenc}
\usepackage[T1]{fontenc}
\usepackage{url}
\usepackage[hidelinks,colorlinks=true,linkcolor=blue]{hyperref}
\usepackage{amsmath,amssymb}
\usepackage{graphicx}
\usepackage{xcolor}
\usepackage{tabularx}
\usepackage{listings}
\usepackage{geometry}
\usepackage{titlesec}
\usepackage{libertine}

\geometry{margin=1in}

\titleformat{\section}{\Large\bfseries\color{darkgray}}{\thesection}{1em}{}
\titleformat{\subsection}{\large\bfseries\color{blue!80!black}}{\thesubsection}{1em}{}

\lstset{
  basicstyle=\small\ttfamily,
  breaklines=true,
  frame=lines,
  backgroundcolor=\color{gray!10}
}

\begin{document}

\begin{center}
    \Huge \textbf{Performance Tuning} \\
    \vspace{0.5em}
    \large \textit{A CEREBRUM Research Publication (v1.2.4)} \\
    \vspace{1em}
    \large \textbf{Bryan Alexander Buchorn} \\
    Independent Researcher \\
    \vspace{2em}
\end{center}

\begin{abstract}
\textbf{Overview:} This guide provides a conceptual entry point into the CEREBRUM framework. It is designed for practitioners and researchers seeking to understand the core reasoning signals of Framework Implementation Guide.
\end{abstract}

\section{CEREBRUM Performance Tuning Guide}

\textbf{Version}: v1.1.0

This guide covers the key parameters affecting query latency, throughput, reasoning quality, and memory usage — and how to tune them for your specific workload.

---

\subsection{The Performance Levers (Quick Reference)}

\begin{center}
{\footnotesize
\begin{tabularx}{\columnwidth}{|>{\raggedright\arraybackslash}X|>{\raggedright\arraybackslash}X|>{\raggedright\arraybackslash}X|>{\raggedright\arraybackslash}X|}
\hline
Parameter & Default & Affects & Tune When \\ \hline
\texttt{beam\_width} & 10 & Latency, recall & Graph is large or queries are slow \\ \hline
\texttt{max\_hop} & 3 & Latency, recall depth & Questions require deeper reasoning \\ \hline
\texttt{n\_trials} (DSCF) & 5 & Community quality, startup time & Startup is slow or Q is low \\ \hline
\texttt{q\_drift\_threshold} & 0.05 & Rebalance frequency & Graph updates are frequent \\ \hline
\texttt{warm\_start\_strength} & 0.0 & First-hop variance & Probab\-ilistic mode on sparse graphs \\ \hline
\texttt{window\_size} (stream) & 10,000 & Memory, temporal context & High-volume ingest \\ \hline
\texttt{batch\_size} (stream) & 50 & Ingest throughput, latency & High-frequency event streams \\ \hline
\texttt{CEREBRUM\_WORKERS} & 4 & Concurrent query throughput & Multi-user deployments \\ \hline
\end{tabularx}
}
\end{center}

---

\subsection{1. Query Latency}

\subsubsection{beam\_width}
The single most impactful parameter for latency. Beam width controls how many candidates are evaluated at each hop.

\begin{center}
{\footnotesize
\begin{tabularx}{\columnwidth}{|>{\raggedright\arraybackslash}X|>{\raggedright\arraybackslash}X|>{\raggedright\arraybackslash}X|>{\raggedright\arraybackslash}X|}
\hline
beam\_width & Latency (21-node) & Latency (10K-node) & H@10 (MetaQA-3hop) \\ \hline
3 & 0.9ms & 2.1ms & 0.271 \\ \hline
5 & 1.2ms & 3.4ms & 0.294 \\ \hline
\textbf{10} & \textbf{1.8ms} & \textbf{6.3ms} & \textbf{0.248} \\ \hline
20 & 3.1ms & 14.8ms & 0.331 \\ \hline
50 & 7.4ms & 48.2ms & 0.339 \\ \hline
\end{tabularx}
}
\end{center}

\textbf{Rule of thumb}: \texttt{beam\_width=10} is the quality/latency sweet spot for most graphs. Increase to 20–50 only for high-stakes queries where latency budget is relaxed (\textgreater 50ms acceptable).

\subsubsection{max\_hop}
Each additional hop multiplies the search space by appro\-ximately \texttt{beam\_width × avg\_degree}.

\begin{lstlisting}
Total candidates ≈ beam_width ^ max_hop
\end{lstlisting}

For a graph with avg\_degree=5 and beam\_width=10:
\begin{itemize}
\item max\_hop=2: ~100 candidates
\item max\_hop=3: ~1,000 candidates
\item max\_hop=4: ~10,000 candidates (consider increasing beam\_width instead)

\end{itemize}
\textbf{Recomm\-endation}: Start at \texttt{max\_hop=3}. Increase to 4 only if 3-hop recall is inadequate on your specific queries.

---

\subsection{2. Community Detection Quality vs. Speed}

\subsubsection{n\_trials (DSCF)}
DSCF is stochastic; running multiple trials and selecting the highest-modularity result improves Q.

\begin{center}
{\footnotesize
\begin{tabularx}{\columnwidth}{|>{\raggedright\arraybackslash}X|>{\raggedright\arraybackslash}X|>{\raggedright\arraybackslash}X|}
\hline
n\_trials & Avg Q (toy\_graph) & Startup time \\ \hline
1 & 0.38 & Fast \\ \hline
3 & 0.41 & Moderate \\ \hline
\textbf{5} & \textbf{0.43} & \textbf{Recommended} \\ \hline
10 & 0.44 & Slow \\ \hline
\end{tabularx}
}
\end{center}

For production, compute communities once at startup or cache them. Communities only need recom\-putation when the graph changes subst\-antially (triggered by \texttt{GlobalRebalancer}).

\subsubsection{Algorithm selection}
\begin{center}
{\footnotesize
\begin{tabularx}{\columnwidth}{|>{\raggedright\arraybackslash}X|>{\raggedright\arraybackslash}X|>{\raggedright\arraybackslash}X|>{\raggedright\arraybackslash}X|}
\hline
Algorithm & Q & Speed & Use when \\ \hline
\texttt{dscf} & Highest & Moderate & Default — best reasoning quality \\ \hline
\texttt{leiden} & High & Moderate & Large graphs where DSCF is slow \\ \hline
\texttt{lpa} & Lower & Fast & Very large graphs (\textgreater 1M nodes) where startup time matters \\ \hline
\end{tabularx}
}
\end{center}

---

\subsection{3. Streaming Ingest Throughput}

\subsubsection{batch\_size and flush\_interval\_ms}
Events are batched before commit to amortize graph write cost.

\begin{lstlisting}
StreamAdapter(
    base_adapter=adapter,
    batch_size=50,              # events per commit — increase for high-volume ingest
    flush_interval_ms=100,      # max wait before commit — decrease for low-latency
    window_size=10_000,         # recency buffer size — reduce to save memory
)
\end{lstlisting}

\begin{center}
{\footnotesize
\begin{tabularx}{\columnwidth}{|>{\raggedright\arraybackslash}X|>{\raggedright\arraybackslash}X|>{\raggedright\arraybackslash}X|>{\raggedright\arraybackslash}X|}
\hline
Scenario & batch\_size & flush\_interval\_ms & Throughput \\ \hline
Low-latency IoT & 10 & 20 & ~500 events/s \\ \hline
General streaming & 50 & 100 & ~2,000 events/s \\ \hline
High-throughput ingest & 200 & 500 & ~8,000 events/s \\ \hline
\end{tabularx}
}
\end{center}

\textbf{Memory}: Each event in the sliding window uses ~200 bytes. \texttt{window\_size=10,000} $\to$ ~2MB.

\subsubsection{GlobalRebalancer tuning}
\begin{lstlisting}
GlobalRebalancer(
    adapter,
    q_drift_threshold=0.05,   # lower = more frequent rebalancing (higher quality, more CPU)
    min_interval_s=30.0,      # minimum seconds between rebalances
)
\end{lstlisting}

Rebalancing is expensive ($O(N \log N)$). For high-ingest deployments, set \texttt{min\_interval\_s=120} to prevent rebalance storms during burst traffic.

---

\subsection{4. Probab\-ilistic Mode (Bayesian Beam Search)}

Probab\-ilistic mode adds Thompson sampling overhead (~15\% latency increase) but improves quality on sparse graphs and provides confidence intervals.

\begin{lstlisting}
BeamTraversal(
    probabilistic=True,
    warm_start_strength=5.0,    # 0.0 = uniform prior (cold), 5.0 = CSA-seeded (warm)
)
\end{lstlisting}

\textbf{warm\_start\_strength} guidance:
\begin{itemize}
\item \texttt{0.0} — standard uniform Beta(1,1) prior (v0.4.0 behavior)
\item \texttt{2.0} — mild warm-start (reduces first-hop variance ~40\%)
\item \texttt{5.0} — strong warm-start (reduces first-hop variance ~85\%, recommended for sparse graphs)
\item \texttt{10.0} — aggressive seeding (can reduce diversity; use only on very sparse graphs)

\end{itemize}
---

\subsection{5. Embedding Choice}

\begin{center}
{\footnotesize
\begin{tabularx}{\columnwidth}{|>{\raggedright\arraybackslash}X|>{\raggedright\arraybackslash}X|>{\raggedright\arraybackslash}X|>{\raggedright\arraybackslash}X|}
\hline
Engine & Semantic quality & Speed & Use when \\ \hline
\texttt{RandomEngine(dim=64)} & None (structural only) & Fastest & Graph topology dominates; no text labels \\ \hline
\texttt{RandomEngine(dim=128)} & None & Fast & Larger graphs needing more embedding capacity \\ \hline
\texttt{SentenceEngine} & High & Slow (GPU-accelerated) & Entity names carry semantic meaning \\ \hline
\texttt{RotatEEngine} (trained) & Relational & Medium (after training) & Rich relation diversity \\ \hline
\end{tabularx}
}
\end{center}

For pure structural reasoning (graphs where entity names are opaque IDs), \texttt{RandomEngine} achieves the same H@10 as \texttt{SentenceEngine} because the community consensus term ($\beta$) dominates over semantic similarity ($\alpha$).

---

\subsection{6. ResourceGovernor Budgets}

\begin{lstlisting}
from core.hardware import ResourceGovernor

governor = ResourceGovernor(
    max_cpu_percent=80,         # cap CPU usage across all queries
    max_memory_mb=2048,         # hard memory limit
    query_timeout_ms=5000,      # abort queries exceeding this
    rem_cpu_budget=0.15,        # fraction of CPU reserved for REM cycle
)
\end{lstlisting}

For latency-critical deployments, set \texttt{rem\_cpu\_budget=0.05} to give more CPU to query processing and reduce REM cycle priority.

---

\subsection{7. Memory Optimi\-zation}

\begin{center}
{\footnotesize
\begin{tabularx}{\columnwidth}{|>{\raggedright\arraybackslash}X|>{\raggedright\arraybackslash}X|>{\raggedright\arraybackslash}X|}
\hline
Component & Memory & Reduction strategy \\ \hline
Embedding matrix & \texttt{N × dim × 4 bytes} & Reduce \texttt{dim} (64$\to$32 halves memory) \\ \hline
Community map & \texttt{N × 8 bytes} & Minimal — unavoidable \\ \hline
Sliding window buffer & \texttt{W × 200 bytes} & Reduce \texttt{window\_size} \\ \hline
Materi\-alized path store & Variable & Call \texttt{adapter.clear\_path\_cache()} \\ \hline
InsightEvent graph & Capped at 10K nodes & Reduce \texttt{meta\_insight\_engine.max\_events} \\ \hline
\end{tabularx}
}
\end{center}

For a 100K-node graph with \texttt{dim=64}: embedding matrix $\approx$ 25MB. For 1M nodes: $\approx$ 250MB.

---

\subsection{8. Benchmark Your Deployment}

\begin{lstlisting}
# Run the synthetic benchmark
python benchmarks/synthetic_eval.py --nodes 1000 --edges 5000 --queries 100

# Run baseline comparison
python benchmarks/baseline_comparison.py --algorithm dscf --beam-width 10

# Profile a specific configuration
python -m cProfile -s cumulative benchmarks/synthetic_eval.py 2>&1 | head -30
\end{lstlisting}

See \texttt{benchmarks/README.md} for full benchmark docum\-entation.

---
\textbf{Copyright © 2026 Bryan Alexander Buchorn. All Rights Reserved.}


\vspace{3em}
\begin{center}
    \small \textit{Project CEREBRUM — March 2026 — Hardened Enterprise Security (v1.2.4)}
\end{center}

\end{document}
